% Baccalauréat malgache, série D, 2013 : énoncé.
\TCDBacAttribution{D}{2013}
\noindent \underline{\textsc{Exercice I}} (5 points) 
      \hspace{\stretch{1}} Les parties A et B sont ind\'ependantes. \hspace{\stretch{1}}
      %
      \begin{description}
            \item[A] On dispose d'un d\'e cubique \`a 6 faces num\'erot\'ees de 1 \`a 6,
                     parfaitement \'equilibr\'e et de deux urnes $U_1$ et $U_2$ contenant des boules
                     indiscernables au toucher de telle sorte que :
                     \begin{description}
                           \item[$U_1$] contient 8 boules blanches et 2 boules noires.
                           \item[$U_2$] contient 3 boules noires et 7 boules rouges.
                     \end{description}
                     %
                     \begin{enumerate}
                           \item L'\'epreuve consiste \`a lancer une fois le d\'e puis \`a 
                                 tirer une boule de $U_1$ ou de $U_2$.
                                 Si le num\'ero 6 appara\^it sur la face sup\'erieure du d\'e, on tire une boule dans $U_1$.
                                 Sinon, on tire une boule dans $U_2$.
                                 
                                 Calculer la probabilit\'e de chacun des \'ev\`enements suivants : 
                                 \begin{description}
                                       \item[$A$ :] << Le chiffre 6 appara\^it et on obtient une boule blanche >>.
                                       \item[$B$ :] << Obtenir une boule noire >>.
                                 \end{description}
                           \item On tire au hasard et simultan\'ement deux boules de $U_2$ et une boule de $U_1$.
                                 On d\'esigne par $X$ la variable al\'eatoire \'egale au nombre de boule(s) noire(s)
                                 obtenue(s).
                                 \begin{enumerate}
                                       \item D\'eterminer l'univers image de $X$.
                                       \item \'Etablir la loi de probabilit\'e de $X$.
                                 \end{enumerate}
                     \end{enumerate}
            \item[B.] Etand donn\'ee une s\'erie statistique \`a deux variables $(X,Y)$ dont la droite de 
                      r\'egression de $Y$ en $X$ est : $y = 0,12 x + 7,88$.
                      Sachant que la moyenne $\bar{X} = 51$ et le coefficient de corr\'elation $r = 0,93$.
                      \begin{enumerate}
                            \item D\'eterminer la moyenne arithm\'etique $\bar{Y}$.
                            \item Peut-on avoir un ajustement lin\'eaire par moindre carr\'es\,?
                                  Expliquer.
                                  D\'eterminer une \'equation de la droite de r\'egression $X$ en $Y$.
                      \end{enumerate}
                      On donnera le r\'esultat \`a $10^{-2}$ pr\`es.
      \end{description}
%%%%%%%%%%%%%%%%%%%%%%%%%%%%%%%%%%%%%%%%%%%%%%%%%%%%%%%%%%%%%%%%%%%%%%%%%%%%%%%%%%%%%%%%%%%%%%%
%%%%%%%%%%%%%%%%%%%                   %%%%%%%%%%%%%%%%%%%                   %%%%%%%%%%%%%%%%%%%
%%%%%%%%%%%%%%%%%%%%%%%%%%%%%%%%%%%%%%%%%%%%%%%%%%%%%%%%%%%%%%%%%%%%%%%%%%%%%%%%%%%%%%%%%%%%%%%
\noindent \underline{\textsc{Exercice II}} (5 points) \newline
      Soit $P(z) = z^3 - (3+3\sqrt{3}i) z^2 - (6-6\sqrt{3}i)z + 8 + 24\sqrt{3} i$ o\`u $z\in\mathbb{C}$.
      \begin{enumerate}
            \item \begin{enumerate}
                        \item D\'emontrer que l'\'equation $P(z) = 0$ admet deux solutions r\'eelles que l'on pr\'ecisera.
                        \item En d\'eduire la r\'esolution de l'\'equation $P(z) = 0$.
                  \end{enumerate}
            
            %
            Dans tout ce qui suite, on consid\`ere dans le plan complexe $\mathscr{P}$ muni d'un rep\`ere
            orthonorm\'e direct $(O;\vec{u},\vec{v})$ d'unit\'e graphique 1~cm, les points $A$, $B$ et $C$ 
            d'affixes respectives $z_A = -2$, $z_B = 4$ et $z_C = 1 + 3\sqrt{3} i$.
            %
            \item \begin{enumerate}
                        \item Placer les points $A$, $B$ et $C$.
                        \item Donner la forme alg\'ebrique du nombre complexe $Z = \frac{z_B - z_A}{z_C - z_A}$ 
                              puis en d\'eduire la nature du triangle $ABC$.
                  \end{enumerate}
            \item Soit $D$ le point tel que $ABCD$ soit un parall\'elogramme.
                  On note par $S$ la similitude plane directe de centre $A$ et qui transforme $B$ en $D$.
                  \begin{enumerate}
                        \item D\'eterminer l'affixe de $D$ et placer ce point.
                        \item D\'eterminer l'expression complexe de $S$ et pr\'eciser ses \'el\'ements caract\'eristiques.
                  \end{enumerate}
      \end{enumerate}
%%%%%%%%%%%%%%%%%%%%%%%%%%%%%%%%%%%%%%%%%%%%%%%%%%%%%%%%%%%%%%%%%%%%%%%%%%%%%%%%%%%%%%%%%%%%%%%
%%%%%%%%%%%%%%%%%%%                   %%%%%%%%%%%%%%%%%%%                   %%%%%%%%%%%%%%%%%%%
%%%%%%%%%%%%%%%%%%%%%%%%%%%%%%%%%%%%%%%%%%%%%%%%%%%%%%%%%%%%%%%%%%%%%%%%%%%%%%%%%%%%%%%%%%%%%%%
\noindent \underline{\textsc{Probleme}} \newline
      On consid\`ere la fonction $f$, \`a variable r\'eelle $x$, d\'efinie par 
      \[
            f(x) = \begin{multiform}
                         (x+1) \expo{x}   \text{ si } x \geq 0 \\
                         \expo{3x} - 3x   \text{ si } x < 0 \;.
                   \end{multiform}
      \]
      On d\'esigne par $(\mathscr{C})$ sa courbe repr\'esentative dans un rep\`ere orthonorm\'e d'origine $O$
      et d'unit\'e graphique 2~cm.
      \begin{enumerate}
            \item Etudier la continuit\'e de $f$ en $x_0 = 0$.
            \item D\'emontrer que $\displaystyle{\lim_{x\rightarrow 0^-} \frac{f(x)-f(0)}{x} = 0}$
                  et $\displaystyle{\lim_{x\rightarrow0^+} \frac{f(x)-f(0)}{x} = 2}$.
                  (On donnera une interpr\'etation graphique de ces r\'esultats.)
            \item \begin{enumerate}
                        \item Etudier la branche infinie de $(\mathscr{C})$ au voisinage de $+\infty$.
                        \item Prouver que la droite $(\mathscr{D})$ d'\'equation $y = -3x$ est une 
                              asymptote oblique pour la courbe $(\mathscr{C})$ au voisinage de $-\infty$.
                  \end{enumerate}
            \item \begin{enumerate}
                        \item Apr\`es avoir \'etudi\'e le sense de variation de $f$, dresse son tableau de variation.
                        \item D\'emontrer que, pour tout $x$ r\'eel n\'egatif, $\expo{3x} - 3x - 1 > 0$.
                  \end{enumerate}
            \item Tracer dans le m\^eme rep\`ere, en pr\'ecisant les deux demi-tangentes, la droite $(\mathscr{D})$ 
                  et la courbe $(\mathscr{C})$.
            \item On appelle $g$ la restriction de la fonction $f$ \`a $]-\infty ; 0[$.
                  \begin{enumerate}
                        \item D\'emontrer que l'\'equation $g(x) = 2$ admet une solution unique $\alpha$.
                        \item On d\'esigne par $g^{-1}$ la fonction r\'eciproque de $g$.
                              Calculer, en fonction de $\alpha$, ${g^{-1}}'(2)$.
                        \item Construire, dans le m\^eme rep\`ere que $(\mathscr{C})$, la courbe $(\mathscr{C}')$
                              repr\'esentant $g^{-1}$.
                        \item Calculer, en fonction de $\alpha$ et en $\mathrm{cm}^2$, l'aire $\mathscr{A}_{\alpha}$
                              du domaine plan limit\'e par : 
                              la courbe $(\mathscr{C})$, la droite $(\mathscr{D})$, les droites d'\'equations 
                              $x = \alpha$ et $x = 0$.
                  \end{enumerate}
      \end{enumerate}













